\section{Definição}
\begin{definition}
    Sejam
    \begin{itemize}
        \item $A\subseteq \mathbb R^n$ aberto;
        \item $p \in A$ um ponto interior de $A$;
        \item $f:A\to \mathbb R^m$ uma função.
    \end{itemize}
    Dizemos que $f$ é \textbf{diferenciável} em $p$ se existe uma transformação linear $T:\mathbb R^n\to \mathbb R^m$ tal que
    \[
        \lim_{h\to 0} \frac{f(p+h)-f(p)-T(h)}{\|h\|} = 0.
    \]
\end{definition}

\begin{lemma}
    Sejam
    \begin{itemize}
        \item $A\subseteq \mathbb R^n$ aberto;
        \item $p \in A$ um ponto interior de $A$;
        \item $f:A\to \mathbb R^m$ uma função diferenciável em $p$.
    \end{itemize}
    Seja $T$ a transformação linear que satisfaz a definição de diferenciabilidade, ou seja, tal que
    \[
        \lim_{h\to 0} \frac{f(p+h)-f(p)-T(h)}{\|h\|} = 0.
    \]
    Então $T$ é a única tal função, e, para cada $i=1,\dots,n$, temos
    \[
        T(e_i) = \frac{\partial f}{\partial x_i}(p).
    \]
\end{lemma}

\begin{proof}
    Seja dado $i \in \{1,\dots,n\}$.
    Temos que:
    \begin{align*}
        \frac{\partial f}{\partial x_i}(p) &= \lim_{t\to 0} \frac{f(p+te_i)-f(p)}{t} \\
        &= \lim_{t\to 0} \frac{f(p+te_i)-f(p)-T(te_i)+T(te_i)}{t}\\
        &= \lim_{t\to 0} \frac{f(p+te_i)-f(p)-T(te_i)}{t} +\frac{tT(e_i)}{t} \\
        &= \lim_{t\to 0} \frac{f(p+te_i)-f(p)-T(te_i)}{t} + T(e_i) \\
        &= \lim_{t\to 0} \frac{f(p+te_i)-f(p)-T(te_i)}{t\|te_i\|}\|te_i\| + T(e_i) \\
        &= \lim_{t\to 0} \frac{f(p+te_i)-f(p)-T(te_i)}{\|te_i\|}\frac{|t|}{t} + T(e_i) \\
        &= 0 + T(e_i) \\
        &= T(e_i).
    \end{align*}
\end{proof}